% =====================================================================
%  Yêu cầu đề tài được giao (TV1) — nội dung theo workspace nhóm
% =====================================================================
\chuongphu{YÊU CẦU ĐỀ TÀI ĐƯỢC GIAO}

\begin{enumerate}[leftmargin=1.2cm,itemsep=3pt]
  \item \textbf{Tên đề tài (Đề tài 10): Mã Hamming (7,4): mô phỏng mã hóa, gây
        nhiễu, phát hiện và sửa lỗi cho chuỗi ký tự.}
  \item \textbf{Môn học:} Đại số tuyến tính (MT1007) -- Học kỳ 1, năm học
        2026--2027.
  \item \textbf{Giảng viên hướng dẫn:} Cô Nguyễn Xuân Mỹ.
  \item \textbf{Nhiệm vụ cụ thể:}
        \begin{enumerate}[label=\arabic{enumi}.\arabic*.,leftmargin=1.2cm]
          \item Nêu cơ sở lý thuyết của mã Hamming (7,4); trình bày ma trận tạo
                mã và ma trận kiểm tra.
          \item Nêu cơ sở lý thuyết giải thích vì sao ma trận tạo mã và ma trận
                kiểm tra có dạng như vậy. Bổ sung: tham khảo mã Hamming mở rộng
                để kiểm tra được trường hợp có $2$ lỗi.
          \item Cho một đoạn mã: nêu rõ cách tạo mã, cách kiểm tra mã có bị lỗi
                hay không và cách sửa lỗi.
          \item Sản phẩm là phần mềm mô phỏng quá trình: một chuỗi ký tự đi vào
                $\rightarrow$ mã hóa thành chuỗi nhị phân $\rightarrow$ gây
                nhiễu $\rightarrow$ phát hiện lỗi $\rightarrow$ sửa lỗi
                $\rightarrow$ xuất ra chuỗi ký tự.
        \end{enumerate}
  \item \textbf{Tài liệu gợi ý:} bài báo ``Ứng dụng mã Hamming trong kiểm soát
        lỗi bộ nhớ'' (Trần Đỗ Hồn Nhiên và cộng sự) và tài liệu do giảng viên
        cung cấp tại
        \href{https://drive.google.com/file/d/1JExRyxYQcP4sspD4K_Ssjy__ME6sZYpr/view?usp=drive_link}{liên kết
        Google Drive}.
  \item \textbf{Yêu cầu chung:}
        \begin{enumerate}[label=\arabic{enumi}.\arabic*.,leftmargin=1.2cm]
          \item Không sao chép dưới bất kỳ hình thức nào; sinh viên được tham
                khảo và tự viết lại theo mức độ hiểu của nhóm.
          \item Không chụp hình từ tài liệu khác dán vào báo cáo; công thức và
                hình ảnh tham khảo từ các nguồn đều phải vẽ và viết lại.
          \item Sinh viên có thể tham khảo nhưng phải tự viết chương trình. Mọi
                sản phẩm phải được giải quyết bằng thuật toán và phần mềm do
                nhóm làm, không lấy sản phẩm từ nguồn khác vào báo cáo.
        \end{enumerate}
\end{enumerate}
